\chapter{Introduction}
\label{c:introduction}

%%%%%%%%%%%%%%%%%%%%%%%%%%%%%%%%%%%%%%%%%%%%%%%%%%%%%%%%%%%%%%%%%%%%%%%%%%%%%%%%
Simulating the interactions within fermionic many-body systems is a central challenge in modern physics.
Fermions, such as electrons, are fundamental building blocks of matter which are completely indistinguishable from each other.
Because of this indistinguishability, swapping two identical fermions introduces a global antisymmetric phase to the system's wavefunction.
By preventing two fermions from occupying the exact same state, this property enforces the Pauli exclusion principle.
To be able to model these properties without tracking spatial coordinates, physicists rely on the second quantization framework and the Fock space.
By using the occupation number representation, we can describe a quantum system by tracking the occupancy of individual spin-orbitals.

However, translating this mathematical model into a computational one presents a severe bottleneck.
Storing quantum states dynamically using static structures like dense arrays scales exponentially and wastes memory on unnecessary zero-amplitude states.
Furthermore, executing quantum chemical simulations on physical quantum hardware can introduce challenges.
Standard methods, such as the Jordan-Wigner transformation, often map fermionic operators to qubits using nonlocal Pauli strings.
When implementing it on actual quantum devices, this mapping requires deep CNOT circuits and Pauli-Z chains.
These factors can make the constructing and applying of full global unitary matrices highly memory prohibitive.

To overcome these limitations, this thesis presents the algorithmic design and implementation of a C++ simulator specifically tailored for second-quantized fermionic operators.
The goal is to bypass physical hardware constraints and memory limits by using a matrix-free direct state evolution.
This approach eliminates the need for Jordan-Wigner strings and allows an efficient simulations of many-body systems while enforcing the antisymmetry required by the Pauli exclusion principle.

